\documentclass[11pt]{article}
\usepackage{amsmath,amssymb}
\usepackage[margin=1in]{geometry}

\title{Clique Tree Construction via AMD Elimination\\
  {\large Implementation notes for \texttt{clique\_ordering\_md.cc}}}
\date{}

\begin{document}
\maketitle

\section{Overview}

The function \texttt{MakeCliqueTreeMinDegreeFromRowSupports} builds a
supernodal clique tree from a set of row supports (the sparsity pattern of
the constraint matrices) and, optionally, the sparsity of a quadratic cost
matrix~$Q$.  The algorithm proceeds in two phases:
\begin{enumerate}
  \item \textbf{Phase~1} (AMD): compute an elimination ordering that
    minimises fill-in, producing for each vertex~$v$ its \emph{later
    set}~$\ell(v)$ (the living neighbours at the moment $v$ is eliminated).
  \item \textbf{Phase~2} (clique tree construction): detect supernodes,
    build the tree, merge small supernodes, reorder, and compute a
    post-order traversal.
\end{enumerate}

Phase~2 is exposed independently as
\texttt{MakeCliqueTreeFromEliminationOrdering} so that callers can supply
their own ordering.

\section{Phase 1: Bitset Approximate Minimum Degree}

\subsection{Graph representation}

Variables are compacted into $\{0,\ldots,n{-}1\}$.  The adjacency graph is
stored as a dense bitset: each row is a packed array of $\lceil n/64\rceil$
64-bit words, so edge queries and clique-fill operations reduce to bitwise
OR/AND on whole words.

The initial edges come from two sources:
\begin{itemize}
  \item \textbf{Row supports.}  Each row support $\{v_1,\ldots,v_k\}$
    induces a clique: every pair $(v_i,v_j)$ gets an edge.
  \item \textbf{$Q$ sparsity.}  If a sparse symmetric matrix~$Q$ is
    provided, its nonzero pattern adds individual pairwise edges.
\end{itemize}

\subsection{Bucket-queue min-degree}

A bucket queue (doubly-linked lists indexed by degree) gives $O(1)$
extraction of the minimum-degree vertex~$v^*$ at each step.  The
elimination step for $v^*$ is:

\begin{enumerate}
  \item Record $\ell(v^*) \gets \mathrm{Adj}(v^*)$ (the living
    neighbours).
  \item \textbf{Fill-in:} make $\ell(v^*)$ a clique by ORing
    each neighbour's row with the \emph{clique mask} (the OR of all
    neighbour bits).  Only nonzero words of the mask are iterated.
  \item \textbf{Removal:} clear $v^*$'s row and remove it from every
    neighbour's adjacency.  Update degrees and bucket positions.
\end{enumerate}

\subsection{Delayed variables}\label{sec:delayed}

In KKT systems the matrix has saddle-point structure
\[
  K = \begin{pmatrix} H & C^\top \\ C & 0 \end{pmatrix},
\]
where $H\succ 0$ on its support.  The dual (equality-constraint) variables
contribute zero to the diagonal; eliminating them first would create a
$0/0$ pivot.  These variables are passed as the
\texttt{delayed\_variables} parameter.

\paragraph{Unified queue with gate check.}
All variables --- primal and delayed --- enter the bucket queue at
initialisation and compete by minimum degree.  Delayed variables are
treated identically to primals except for a \emph{gate check} applied
when they reach the front of the queue.

When the minimum-degree vertex~$v^*$ is popped:
\begin{itemize}
  \item If $v^*$ is primal, it is eliminated immediately (standard AMD).
  \item If $v^*$ is delayed, it must first \emph{claim} an unclaimed,
    already-eliminated primal from its original (pre-fill-in) neighbours.
    Each primal can be claimed by at most one dual.  If a claimable primal
    is available, $v^*$ is eliminated.  If not, $v^*$ is \emph{skipped}:
    it remains in the queue at its current degree and the next vertex is
    tried.
\end{itemize}

Skipped duals are retried on the next step.  As more primals are
eliminated, more claims become available, so skipped duals eventually
pass.  If no vertex in the queue passes (all remaining vertices are
gated duals with no claimable primal), the first skipped dual is accepted
unconditionally as a fallback.

\subsection{Structural independence guarantee}\label{sec:struct-indep}

Two delayed variables are \emph{structurally dependent} in a supernodal
block if, restricted to the primals within that supernode, their rows of
the constraint matrix~$C$ have identical sparsity patterns.  The local
block is then rank-deficient regardless of the numerical values.

\paragraph{Why the gate prevents this.}
The primal-claim gate ensures that within any supernode, each dual
variable couples to at least one primal that no other dual in the tree
has claimed.  Because the gate requires the claimed primal to be
already eliminated (\texttt{deg[p] < 0}), and the min-degree ordering
naturally places the claimed primal nearby in the elimination order,
the claimed primal will typically be in the same supernode.

For well-posed problems (number of duals~$\le$~number of primals),
every dual can find a unique primal to claim.  Two duals in the same
supernode therefore couple to at least two distinct primals, making
their rows linearly independent.

\paragraph{Original primal adjacency.}
Before the elimination loop begins, for each delayed variable~$v$ we
record $\mathrm{OrigPrimalAdj}(v)$: the sorted list of~$v$'s non-delayed
neighbours in the original (unfilled) graph.  The gate check searches
this list for an unclaimed eliminated primal.  The original adjacency is
used because it reflects the true $C$-matrix coupling; fill-in edges
do not contribute matrix entries.

\section{Phase 2: Clique tree from an elimination ordering}

Phase~2 is implemented in \texttt{MakeCliqueTreeFromEliminationOrdering}.
Its input is an \texttt{EliminationOrdering} containing:
\begin{itemize}
  \item \texttt{order}: the elimination sequence in compact variable
    indices,
  \item $\ell(v)$ (\texttt{later}): the living neighbours of $v$ when
    it was eliminated,
  \item \texttt{parent\_col}$[v]$: the first element of $\ell(v)$ in
    elimination order,
  \item \texttt{unique\_vars}: the compact $\to$ original variable
    mapping.
\end{itemize}

\subsection{Later-set sorting and parent computation}

For each vertex in elimination order, $\ell(v)$ is sorted by elimination
position.  The first element becomes $\mathrm{parent\_col}(v)$: the
variable eliminated immediately after~$v$ among its neighbours.  This
defines the \emph{elimination tree}.

\subsection{Supernode detection}

A \emph{supernode} is a maximal chain
$v_1, v_2, \ldots, v_s$ of consecutively eliminated vertices such that for
each consecutive pair $(v_i, v_{i+1})$:
\begin{enumerate}
  \item $v_{i+1} = \mathrm{parent\_col}(v_i)$ (the parent in the
    elimination tree),
  \item $v_{i+1}$ is eliminated immediately after~$v_i$
    ($\mathrm{pos}(v_{i+1}) = \mathrm{pos}(v_i) + 1$),
  \item $\ell(v_i) = \{v_{i+1}\} \cup \ell(v_{i+1})$ (identical fill
    pattern shifted by one).
\end{enumerate}

Condition~3 means that eliminating $v_i$ creates no new fill beyond the
edge to $v_{i+1}$; the two columns can be combined into one dense
block.  The algorithm scans the elimination order and greedily extends
each chain until one of the three conditions fails.

\subsection{Clique and tree construction}

For each supernode~$S_i = \{v_1,\ldots,v_s\}$, the \emph{maximal clique}
is
\[
  C_i = S_i \cup \ell(v_1),
\]
that is, the supernode columns plus all of their shared later neighbours.
Variables are mapped back to the original numbering via
\texttt{unique\_vars}.

The \emph{tree parent} of supernode~$S_i$ is the supernode containing
$\mathrm{parent\_col}(v_s)$, the parent of the last column in the chain.

The \emph{separator} of $S_i$ is $C_i \cap C_{\mathrm{parent}(i)}$: the
variables shared with the parent's clique.  The supernode is
$C_i \setminus \text{separator}$: the variables owned exclusively by this
node.

\subsection{Supernode merging}

Small supernodes (size $\le$ \texttt{max\_merge\_supernode\_size}) are
merged into their parents.  Merging absorbs the child's supernode variables
into the parent's supernode.  Children of the merged node are reparented to
the grandparent.  This is iterated until no more merges are possible, then
the tree is compacted to remove empty nodes.

Merging reduces the number of tree nodes (and thus the number of
factorisation steps), at the cost of potentially enlarging the dense blocks.

\subsection{Supernode reordering}

Within each clique the supernode variables are reordered using one of:
\begin{itemize}
  \item \texttt{SUPERNODE\_REORDER\_BFS\_GREEDY}: BFS-based greedy
    ordering that tries to place variables contiguously.
  \item \texttt{SUPERNODE\_REORDER\_PQ\_TREE}: PQ-tree based consecutive
    arrangement.
  \item \texttt{SUPERNODE\_REORDER\_NONE}: keep the ordering from the
    elimination phase.
\end{itemize}

\subsection{Post-order traversal}

Finally, a post-order traversal of the tree is computed via iterative DFS.
The result is stored in
\texttt{ct.post\_order\_position\_to\_clique}, which gives the node index
for each position in the traversal.  The supernodal Cholesky factorisation
processes nodes in this order so that children are factored before their
parents.

\section{Output}

The returned \texttt{CliqueTree} contains:
\begin{itemize}
  \item \texttt{supernodes}$[i]$: variables owned by node~$i$ (in
    original numbering),
  \item \texttt{separators}$[i]$: variables shared with the parent,
  \item \texttt{node\_to\_parent}$[i]$: parent node index ($-1$ for
    roots),
  \item \texttt{post\_order\_position\_to\_clique}: post-order traversal.
\end{itemize}

The optional output \texttt{maximal\_cliques\_out} contains
$\texttt{supernodes}[i] \cup \texttt{separators}[i]$ for each node,
corresponding to the maximal cliques of the filled graph.  Constraint
decomposition (\texttt{Decompose}) uses these to assign each constraint
row to the smallest containing clique.

\end{document}
